\subsection{Локальные экстремумы функций. Теоремы Ферма и Ролля.}

\subsubsection*{Определения экстремумов}

Пусть функция $f(x)$ определена на интервале $(a, b)$.
Пусть $x_0 \in (a, b)$.

\begin{tcolorbox}[title=Определение (Локальный экстремум)]
    Точка $x_0$ называется точкой \textbf{максимума} (или минимума), если существует окрестность $U(x_0) \subset (a, b)$, такая что для всех $x \in U(x_0)$:
    \[ f(x_0) \ge f(x) \quad (\text{максимум}) \]
    \[ f(x_0) \le f(x) \quad (\text{минимум}) \]
\end{tcolorbox}

\begin{tcolorbox}[title=Определение (Строгий экстремум)]
    Точка $x_0$ называется точкой \textbf{строгого максимума} (или минимума), если для всех $x$ из проколотой окрестности $\mathring{U}(x_0)$ (т.е. $x \neq x_0$):
    \[ f(x) < f(x_0) \quad (\text{строгий максимум}) \]
    \[ f(x) > f(x_0) \quad (\text{строгий минимум}) \]
\end{tcolorbox}

Точки максимума и минимума объединяются общим названием — \textbf{точки экстремума}.

\subsubsection*{Теорема Ферма}

\begin{tcolorbox}[title=Теорема Ферма (Необходимое условие экстремума)]
    Пусть функция $f(x)$ определена на интервале $(a, b)$ и $x_0 \in (a, b)$ — точка экстремума.
    Если функция $f(x)$ дифференцируема в точке $x_0$, то её производная в этой точке равна нулю:
    \[ f'(x_0) = 0 \]
\end{tcolorbox}

\textbf{Замечание:} В точке экстремума производная либо равна нулю, либо не существует.

\textbf{Доказательство:}
По условию $f(x)$ дифференцируема в $x_0$. Запишем определение дифференцируемости (через приращение):
\[ f(x) - f(x_0) = f'(x_0)(x - x_0) + \alpha(x)(x - x_0), \]
где $\alpha(x) \to 0$ при $x \to x_0$.
Вынесем $(x - x_0)$ за скобку:
\[ f(x) - f(x_0) = (x - x_0) \left[ f'(x_0) + \alpha(x) \right] \]

Предположим противное: пусть $f'(x_0) \neq 0$.
Так как $\alpha(x) \to 0$, то в некоторой малой окрестности $U(x_0)$ выражение в скобках $\left[ f'(x_0) + \alpha(x) \right]$ сохраняет знак (тот же, что и у $f'(x_0)$).
Однако множитель $(x - x_0)$ меняет знак при переходе через точку $x_0$ (слева минус, справа плюс).
Следовательно, всё произведение (а значит и разность $f(x) - f(x_0)$) меняет знак в окрестности точки $x_0$.

Это означает, что $f(x) - f(x_0)$ принимает как положительные, так и отрицательные значения, что противоречит определению экстремума (где разность должна быть либо всегда $\le 0$, либо всегда $\ge 0$).
Противоречие $\Rightarrow f'(x_0) = 0$.

\subsubsection*{Теорема Ролля}

\begin{tcolorbox}[title=Теорема Ролля (О нулях производной)]
    Пусть функция $f(x)$:
    \begin{enumerate}
        \item Непрерывна на отрезке $[a, b]$ ($f \in C[a, b]$).
        \item Дифференцируема на интервале $(a, b)$.
        \item Принимает равные значения на концах: $f(a) = f(b)$.
    \end{enumerate}
    Тогда существует хотя бы одна точка $\xi \in (a, b)$, такая что:
    \[ f'(\xi) = 0 \]
\end{tcolorbox}

\textbf{Доказательство:}
Так как $f(x)$ непрерывна на отрезке $[a, b]$, то по \textbf{второй теореме Вейерштрасса} она достигает на этом отрезке своего глобального максимума ($M$) и минимума ($m$).
То есть существуют точки $x_{max}$ и $x_{min}$ такие, что $f(x_{max}) = M, f(x_{min}) = m$.

Рассмотрим два случая:
\begin{enumerate}
    \item Пусть $M = m$.
    Тогда $f(x) = \text{const}$ на всём отрезке. Производная константы равна 0 в любой точке интервала. Теорема верна.

    \item Пусть $M > m$.
    Так как по условию $f(a) = f(b)$, то значения $M$ и $m$ не могут одновременно достигаться на концах отрезка (иначе было бы $M=m$).
    Значит, хотя бы одна из точек экстремума (обозначим её $\xi$) находится строго внутри интервала: $\xi \in (a, b)$.
    
    Так как $\xi$ — точка экстремума и $\xi \in (a, b)$, и функция дифференцируема в этой точке, то по \textbf{теореме Ферма}:
    \[ f'(\xi) = 0. \]
\end{enumerate}
Теорема доказана.


